\section*{Chapitre \ref{ch:iv3:operating-systems-concurrency-and-networks} --- Systèmes d'exploitation, concurrence et réseaux}
\markboth{\protect\omTocMark{}{Réponses~: Systèmes d'exploitation, concurrence et réseaux}}{\protect\omTocMark{}{Réponses~: Systèmes d'exploitation, concurrence et réseaux}}

\begin{solution}{iq:iv3:operating-systems-concurrency-and-networks:1}
Il affiche \texttt{after fork 0} et \texttt{after thread 1}. L'enfant possède sa propre copie de \texttt{count}, donc son incrément reste dans l'enfant, tandis que le thread modifie l'unique \texttt{count} du processus.

Les threads d'un processus partagent l'espace d'adressage, donc ils partagent les données globales et le tas. Ils partagent aussi les fichiers ouverts et les gestionnaires de signaux. Chaque thread possède sa propre pile, ses propres registres et son propre compteur de programme, ainsi que son propre stockage local au thread. Un thread peut néanmoins atteindre la pile d'un autre au moyen d'un pointeur, puisque la mémoire ne forme qu'un seul espace.

Après \texttt{fork}, l'enfant démarre avec une copie de la mémoire du parent, que Linux réalise de façon paresseuse, et avec une copie de sa table des descripteurs de fichiers. Aucune mémoire ne reste partagée, sauf si elle a été mappée comme partagée. Les descriptions de fichiers ouverts font exception~: le parent et l'enfant partagent les positions de lecture-écriture par leur intermédiaire.

\iqwrong{«~Un processus et un thread ne diffèrent que par leur coût.~» Le coût découle de la vraie différence, qui est de savoir si les écritures sont visibles de l'autre côté.}
\iqfollow{Lequel est le plus sûr lorsqu'un des deux plante~?}{Processus~: une écriture sauvage dans l'un ne peut pas corrompre la mémoire de l'autre, et un plantage ne tue que ce processus. Une faute de segmentation dans un thread met fin à tous les autres.}
\end{solution}

\begin{solution}{iq:iv3:operating-systems-concurrency-and-networks:2}
Il affiche \texttt{4000 1000 0}. \texttt{shared} est une variable atomique unique que les quatre threads incrémentent. \texttt{mine} est une variable par thread, si bien que chaque thread compte jusqu'à 1000, et le \texttt{mine} du thread principal, qui n'a fait aucun comptage, vaut toujours 0.

\texttt{shared} possède une seule adresse, dans le segment de données du programme. Chaque thread reçoit son propre bloc de stockage local au thread au moment de son démarrage, et \texttt{mine} s'y trouve à un décalage fixe. Sous Linux x86-64, un registre (\texttt{fs}) pointe vers le bloc du thread courant, si bien qu'accéder à \texttt{mine} ne demande ni verrou ni atomique.

Cela convient à une statistique mise à jour souvent et lue rarement. Chaque thread incrémente un compteur \texttt{thread\_local} qui lui est propre, qu'aucun autre cœur n'écrit, et un lecteur en fait la somme de temps à autre. On évite ainsi qu'une ligne de cache circule d'un cœur à l'autre à chaque incrément, au prix d'une somme qui n'est pas un instantané, et les compteurs doivent être enregistrés pour qu'un lecteur puisse les retrouver.

\iqwrong{«~\texttt{thread\_local} est identique à une variable locale.~» Une variable locale est recréée à chaque appel, alors que celle-ci vit aussi longtemps que son thread et conserve sa valeur entre les appels.}
\end{solution}

\begin{solution}{iq:iv3:operating-systems-concurrency-and-networks:3}
Il affiche \texttt{state Z}, \texttt{code 7} et \texttt{gone True}. Un zombie est un enfant qui s'est terminé sans avoir encore été attendu. Il ne détient ni mémoire ni fichiers ouverts, seulement son entrée dans la table des processus, qui conserve le statut de sortie pour le parent.

Pour l'éliminer, il faut que le parent appelle \texttt{wait} ou \texttt{waitpid}, comme le fait le listing, et cela donne au parent le code 7. On ne peut pas tuer un zombie, puisqu'il est déjà mort, si bien que \texttt{kill -9} ne fait rien.

Si le parent meurt le premier, ses enfants sont confiés à \texttt{init} ou au plus proche subreaper, et ce processus les attend, de sorte que les zombies disparaissent. Le problème vient d'un parent de longue durée qui n'attend jamais~: chaque zombie occupe un identifiant de processus, et le système n'en dispose que d'un nombre limité.

Un parent qui ne se soucie pas du statut peut passer \texttt{SIGCHLD} à «~ignoré~», et alors aucun zombie ne subsiste.

\iqwrong{«~Un zombie est un processus bloqué qui consomme du CPU.~» Il ne consomme aucun CPU, puisqu'il s'est déjà terminé.}
\end{solution}

\begin{solution}{iq:iv3:operating-systems-concurrency-and-networks:4}
Il affiche \texttt{b\textquotesingle{}abc\textquotesingle{} b\textquotesingle{}de\textquotesingle{}} et \texttt{b\textquotesingle{}abcde\textquotesingle{}}. Le nombre d'appels \texttt{recv} pour TCP n'est pas fixe~: deux ou trois ici, en comptant le dernier qui ne renvoie rien à la fermeture, et cela dépend du calendrier et des tampons. UDP préserve les frontières des messages, alors que TCP est un flux d'octets sans frontières.

TCP garantit que les octets arrivent complets, dans l'ordre et sans doublons, ou que la connexion signale un échec. Il ajoute la retransmission, le contrôle de flux et le contrôle de congestion, et il exige d'abord l'établissement d'une connexion. UDP envoie chaque datagramme de manière indépendante, qui peut être perdu, dupliqué ou réordonné, et il n'informe pas l'émetteur de quoi que ce soit. Un datagramme qui arrive arrive entier.

L'angle du trading est le coût de l'ordre. Un segment TCP perdu retient chaque octet suivant jusqu'à sa réémission, alors qu'un consommateur de données de marché préfère disposer tout de suite des prix les plus récents. Les flux de données sont donc généralement envoyés en UDP avec des numéros de séquence, et les trous sont comblés par le consommateur lui-même, tandis que la saisie d'ordres passe en général par TCP, puisque chaque message doit arriver et dans l'ordre.

\iqwrong{«~UDP n'est pas fiable, il perd des données en permanence.~» Sur un réseau sain, cela arrive rarement. Le point essentiel est que rien ne le répare lorsque cela se produit.}
\end{solution}

\begin{solution}{iq:iv3:operating-systems-concurrency-and-networks:5}
Il affiche \texttt{little}, \texttt{04030201 01020304} et \texttt{0x4030201}. x86-64 range d'abord l'octet de poids faible d'un nombre, si bien que 0x01020304 est disposé comme 04 03 02 01. L'ordre des octets du réseau est big-endian, c'est-à-dire que l'octet de poids fort vient en premier~; c'est ce que produit le format \texttt{!I} et ce vers quoi convertit \texttt{htonl}.

Ainsi, \texttt{htonl} inverse les octets sur cette machine. Sur une machine big-endian, il ne ferait rien.

Un flux qui envoie ses prix en ordre little-endian se lit directement avec le format natif sur x86-64, sans inversion. Sur une machine big-endian, le même prix exige un format little-endian explicite. Pour du code portable, il faut écrire l'ordre voulu~: \texttt{<I} signifie little-endian sur toute machine, alors que \texttt{=I} désigne ce qu'utilise la machine.

\iqwrong{«~L'ordre des octets du réseau est ce qu'utilise la carte réseau, donc le noyau s'en charge.~» Le noyau convertit ses propres en-têtes. Les champs de votre charge utile vous reviennent.}
\iqfollow{À quoi ressemble un mauvais ordre des octets~?}{Une petite valeur devient énorme et une grande devient petite, par un facteur pouvant atteindre $2^{24}$ (exactement ce facteur pour une valeur qui tient sur un octet), ce qui doit inciter à vérifier tout prix.}
\end{solution}

\begin{solution}{iq:iv3:operating-systems-concurrency-and-networks:6}
Il affiche \texttt{True}, \texttt{False} et \texttt{b\textquotesingle{}a\textquotesingle{} b\textquotesingle{}b\textquotesingle{}}. Les deux connexions ont la même adresse et le même port locaux, et des ports distants différents. Une connexion TCP est identifiée par quatre éléments~: adresse locale, port local, adresse distante et port distant. Avec le protocole, cela fait le quintuplet.

Un seul port serveur peut donc porter de nombreuses connexions. Le noyau recherche chaque segment entrant par ce quintuplet et le remet au bon socket. \texttt{accept} renvoie un nouveau socket pour chaque connexion, et le socket d'écoute continue d'écouter.

Les véritables limites sont ailleurs. Un client qui communique avec une seule adresse et un seul port serveur ne dispose que de sa plage de ports source éphémères, de quelques dizaines de milliers ici (\texttt{ip\_local\_port\_range}), et la limite d'un serveur vient de la mémoire et des descripteurs de fichiers.

\iqwrong{«~Chaque client a besoin de son propre port serveur.~» Cela confond le port d'écoute avec la connexion. Aucun nouveau port n'est ouvert lors de \texttt{accept}.}
\iqfollow{En quoi un socket UDP est-il différent~?}{Il n'a pas de connexion, si bien qu'un seul socket reçoit de tous les émetteurs, et l'adresse source accompagne chaque datagramme.}
\end{solution}

\begin{solution}{iq:iv3:operating-systems-concurrency-and-networks:7}
Il affiche \texttt{1110} et \texttt{10}. Trois emplacements peuvent être pris et la quatrième tentative échoue. Après la libération par l'autre thread, un emplacement est de nouveau libre. Un sémaphore est un compteur sans propriétaire, si bien que n'importe quel thread peut le libérer, alors qu'un mutex a un propriétaire, et seul ce thread peut le déverrouiller.

Un sémaphore convient donc au comptage~: au plus $N$ requêtes en cours, ou des éléments disponibles dans une file, un thread acquérant et un autre libérant. Il sert aussi à la signalisation entre threads, ce qu'un mutex ne peut pas faire. Un mutex protège une section critique, et sa règle de propriété permet à un système d'offrir l'héritage de priorité ou une vérification de débogage.

Déverrouiller un \texttt{std::mutex} depuis un thread qui ne le possède pas constitue un comportement indéfini. Un sémaphore binaire ressemble à un mutex, mais il n'offre pas cette protection, puisque rien ne lie la libération à celui qui a acquis.

\iqwrong{«~Un sémaphore binaire et un mutex sont la même chose.~» Les deux ont les mêmes états, mais seul un mutex a un propriétaire.}
\iqfollow{\texttt{try\_acquire} peut-il échouer alors qu'un emplacement est libre~?}{La norme autorise un échec parasite, si bien qu'un programme réel ne doit pas tenir l'échec pour une preuve qu'aucun emplacement n'est libre. Cela ne s'est pas produit dans nos exécutions, qui n'avaient aucune contention.}
\end{solution}

\begin{solution}{iq:iv3:operating-systems-concurrency-and-networks:8}
Il affiche \texttt{caught [15]}, \texttt{refused} et \texttt{-9}. \texttt{SIGTERM} est une demande, que le processus peut traiter. \texttt{SIGKILL} ne peut être ni interceptée, ni bloquée, ni ignorée, car l'opérateur a besoin d'un moyen qui fonctionne toujours. Le statut $-9$ montre que l'enfant a été tué par ce signal.

Le noyau met simplement fin à un processus sur \texttt{SIGKILL}. Aucun gestionnaire ne s'exécute, aucun destructeur n'est appelé et aucune donnée tamponnée n'est vidée.

Un processus de trading doit donc recevoir d'abord \texttt{SIGTERM}, et son gestionnaire doit lever un drapeau puis rendre la main. La boucle principale cesse alors d'envoyer de nouveaux ordres, annule ceux qui sont ouverts, vide son journal et se termine. Ce n'est que s'il ne parvient pas à finir à temps que l'opérateur envoie \texttt{SIGKILL}. Et comme le processus peut mourir malgré tout, par exemple lors d'un plantage, l'annulation à la déconnexion (cancel on disconnect) côté plateforme, lorsqu'une plateforme la propose, est la protection qui ne dépend pas du processus.

Un gestionnaire doit faire très peu de choses. La plupart des appels de bibliothèque ne sont pas sûrs à l'intérieur d'un tel gestionnaire, car le signal peut arriver alors que cette même bibliothèque détient un verrou.

\iqwrong{«~Mettre la logique d'annulation dans le gestionnaire de \texttt{SIGKILL}.~» Il n'existe pas de tel gestionnaire. Elle doit se trouver sur \texttt{SIGTERM}, et un plantage doit être couvert ailleurs.}
\end{solution}

\begin{solution}{iq:iv3:operating-systems-concurrency-and-networks:9}
L'enfant subit environ 0 défauts à la lecture et environ 2000 à l'écriture. Dans l'exécution présentée ici, les deux comptes valaient exactement 0 et 2000. L'enfant affiche 2 et le parent 1, à la même adresse, puisque chaque processus possède ses propres tables de pages qui associent une adresse virtuelle à des pages physiques différentes.

\texttt{fork} copie les tables de pages et passe les pages inscriptibles en lecture seule dans les deux processus. Une lecture tombe donc sur une page déjà mappée, et aucun défaut ne se produit. La première écriture dans chaque page provoque un défaut. Le noyau copie alors cette page, donne au scripteur sa propre copie et la rend inscriptible. C'est la copie sur écriture. Ainsi, \texttt{fork} lui-même est peu coûteux en données, et le coût survient plus tard, page par page.

Pour un grand processus qui forke afin de sauvegarder un instantané, c'est une bonne astuce~: l'enfant voit une vue figée et l'écrit, tandis que le parent continue de tourner. Mais chaque page que le parent modifie pendant la sauvegarde coûte désormais un défaut, ainsi qu'une copie de 4 Kio. Si le parent réécrit la plus grande partie de sa mémoire, cela double au pire l'usage mémoire, et cela ajoute de la latence aux écritures du chemin critique, réparties de façon imprévisible dans le temps. \texttt{fork} lui-même doit aussi copier les tables de pages, ce qui croît avec la taille mappée, et pendant ce temps le processus est arrêté.

Le listing désactive les grandes pages pour cette plage, afin que le premier accès du parent ne dépende pas du réglage du système. Sur un noyau récent, les comptes de l'enfant ne changent pas~: une écriture dans une grande page partagée scinde le mappage et ne copie que 4 Kio, et une exécution ici avec les grandes pages activées a de nouveau donné 0 et 2000. Les anciens noyaux tentaient de copier la grande page entière.

\iqwrong{«~L'enfant reçoit une copie de toute la mémoire au moment de \texttt{fork}.~» Il reçoit le mappage. Les copies physiques apparaissent à la première écriture.}
\end{solution}

\begin{solution}{iq:iv3:operating-systems-concurrency-and-networks:10}
Les indices valent 254, 72, 418 et 359, avec un décalage de 2748 (0xabc)~: bits 47--39, 38--30, 29--21, 20--12 et 11--0. Mapper 1 Gio demande 262\,144 entrées, soit 512 tables de 4 Kio, c'est-à-dire 2 Mio, plus 3 tables supérieures~: 2\,109\,440 octets au total. Avec des pages de 2 Mio, le dernier niveau disparaît, et il faut 3 tables, soit 12 Kio.
\omcode{code/interviews-3/05-operating-systems-concurrency-and-networks/python/iv_page_walk.py}{1}{17}{Découpage d'une adresse, et comptage des tables.\protect\hfil\protect\break}\label{lst:iv3:os:pagewalk}

Chaque table compte 512 entrées, soit $2^9$, si bien que chaque niveau occupe 9 bits. Une page de 4096 octets occupe 12 bits pour le décalage. $4 \times 9 + 12 = 48$. Sous Linux, l'espace utilisateur occupe la moitié basse, sous 47 bits, soit environ 128 To.

La boucle remonte les niveaux. $2^{30} / 2^{12} = 262\,144$ pages demandent autant d'entrées au dernier niveau, qui remplissent 512 tables~; celles-ci demandent 512 entrées un niveau au-dessus, dans 1 table~; et ainsi de suite jusqu'au sommet. Le surcoût vaut donc environ 0.2~\% de ce qui est mappé. Le test vérifie la découpe en reconstruisant 1000 adresses aléatoires à partir de leurs morceaux, et le nombre de tables contre une marche lente sur chaque page, qui recueille les tables distinctes. Une seule page nécessite tout de même une table à chaque niveau.

Ce ne sont pas les tables de pages qui rendent les grandes pages attrayantes. Une entrée du cache de traduction couvre alors 2 Mio et non 4 Kio, si bien que les échecs de cache sont beaucoup moins nombreux, et chaque échec a un niveau de moins à parcourir.

\iqwrong{«~Les tables de pages d'un processus occupent 8 octets pour chaque octet.~» Elles occupent 8 octets pour chaque page, ce qui explique que le surcoût soit faible.}
\iqfollow{Et qu'en est-il de la mémoire mappée mais jamais touchée~?}{Linux ne remplit les tables que pour les pages accédées, si bien qu'un mappage inactif ne coûte presque rien.}
\end{solution}

\begin{solution}{iq:iv3:operating-systems-concurrency-and-networks:11}
Le programme affiche \texttt{heap mmap}, \texttt{unmapped} et \texttt{heap}. Les petits blocs viennent du tas, découpé par \texttt{malloc} lui-même. Un bloc au-delà d'un seuil, initialement de 128 Kio, obtient son propre mappage auprès du noyau, que \texttt{free} rend aussitôt~; la deuxième ligne le trouve donc non mappé.

La dernière ligne surprend souvent. Le seuil de glibc bouge. Lorsqu'un bloc situé au-dessus du seuil et en deçà de sa limite haute est libéré, le seuil monte à la taille de ce bloc. Ainsi, après la libération de 1 Mio, la requête suivante de 1 Mio vient du tas. L'idée est qu'un programme qui libère et réalloue de gros blocs ne doit pas payer un appel système à chaque fois.

Pour la latence, cela compte de trois façons. Un \texttt{mmap} et un \texttt{munmap} sont des appels système. Un mappage neuf est vide, si bien que chaque nouvelle page provoque un défaut de page au premier accès. Et comme le seuil bouge, une même requête peut coûter différemment selon l'historique~; un benchmark peut donc ne pas dire ce que fait la première exécution.

Sur le chemin critique, le remède consiste à ne rien allouer là. Il faut allouer dès le départ, toucher les pages et réutiliser la mémoire. Pour supprimer le seuil mobile, il faut fixer \texttt{M\_MMAP\_THRESHOLD} avec \texttt{mallopt}, ce qui désactive le comportement dynamique.

\iqwrong{«~\texttt{free} rend la mémoire au système d'exploitation.~» Seulement pour un bloc mappé, ou pour le sommet du tas au-delà d'un seuil de réduction. La plupart des libérations la rendent simplement aux listes de \texttt{malloc}.}
\end{solution}

\begin{solution}{iq:iv3:operating-systems-concurrency-and-networks:12}
Le programme affiche environ \texttt{1024 0}, puis \texttt{0 100}. L'appel ne donne que de l'espace d'adressage, sans mémoire physique derrière. Une page reçoit de la mémoire physique à son premier accès, par un défaut de page. La taille virtuelle saute donc de 1024 d'un coup, et la taille résidente ne grandit que pour les 100 Mio effectivement accédés.

La première ligne est une promesse, et la seconde la réalité. Le noyau s'en sert pour l'allocation paresseuse, pour la copie sur écriture de fork, et pour les programmes qui mappent bien plus qu'ils n'utilisent. Dans son mode par défaut, il ne refuse que les surengagements manifestes~; une demande de 1 Gio ici réussit donc, tandis qu'une demande bien supérieure à ce que la machine peut fournir échouerait. La règle exacte dépend d'un paramètre du système.

Si chaque processus accède à tout ce qu'il a demandé et que la mémoire vient à manquer, l'échec ne vient pas de \texttt{mmap}, qui a rendu la main depuis longtemps. Il survient à un accès ultérieur, où le noyau n'a plus de pages à donner~; il désigne alors une victime avec le out-of-memory killer. Un processus peut donc mourir à cause d'une simple écriture banale, sur une ligne sans aucune allocation.

Pour un processus de trading, cela suggère d'allouer et de toucher la mémoire au démarrage, et de la verrouiller avec \texttt{mlock} si un défaut de page sur le chemin critique est inacceptable. Un échec au démarrage apparaît alors avant l'ouverture du marché.

\iqwrong{«~Vérifier la valeur de retour de \texttt{malloc} protège de l'épuisement de la mémoire.~» Cela aide contre les limites d'espace d'adressage. Sur un système qui surengage, un pointeur non nul ne garantit aucune page.}
\end{solution}

\begin{solution}{iq:iv3:operating-systems-concurrency-and-networks:13}
Avec 1000, le programme affiche 1000 et utilise environ 1 Mio. 100\,000 requiert environ 100 Mio sous une limite de 8 Mio, et le processus meurt avec \texttt{SIGSEGV}. Sous la limite de 256 Mio, il affiche 100000. On ne peut pas intercepter cela avec \texttt{try}~: un débordement de pile est une faute de mémoire, et non une exception C++. La pile d'un nouveau thread fait 8 Mio sous cette limite, ce que relit notre test.

La pile croît vers le bas page par page, au fur et à mesure des accès. Lorsqu'elle atteint la limite fixée par \texttt{ulimit -s}, le noyau refuse de l'étendre, et l'accès provoque une faute d'accès. Un gestionnaire de signal pour ce cas aurait lui-même besoin d'une pile~; il doit donc s'exécuter sur une pile de secours (\texttt{sigaltstack}), et même alors il peut peu faire d'autre que signaler l'erreur.

Un nouveau thread ne croît pas de cette façon. Toute sa pile est mappée à sa création, avec une taille fixée par \texttt{RLIMIT\_STACK} au démarrage du programme. Si cette limite est illimitée, elle vaut 2 Mio sur x86-64, si bien que modifier \texttt{ulimit} peut changer la taille de la pile de chaque thread de manière surprenante. Il faut la fixer explicitement avec \texttt{pthread\_attr\_setstacksize} lorsque cela importe.

Ce mappage est de la mémoire virtuelle~; 1000 threads inactifs de 8 Mio occupent donc 8 Go d'espace d'adressage, et peu de mémoire réelle. Des pages de garde se trouvent sous chaque pile, ce qui transforme la plupart des débordements en faute, sauf si un cadre de pile de grande taille saute par-dessus la garde.

\iqwrong{«~La profondeur de récursion est limitée par le langage.~» Le C++ n'a pas de telle limite. La pile et son réglage en imposent une, qui varie d'un thread à l'autre et d'une machine à l'autre.}
\end{solution}

\begin{solution}{iq:iv3:operating-systems-concurrency-and-networks:14}
La plus petite valeur possible est 1000~: sur un seul cœur, l'enfant doit s'exécuter entre l'écriture du parent et le retour de la lecture du parent, si bien que le parent doit être retiré du cœur au moins une fois à chaque tour. Le programme a affiché 1000 ici. Un changement de contexte sauvegarde les registres d'un thread, en charge ceux d'un autre, ainsi que son espace d'adressage s'il appartient à un autre processus.

Le listing distingue les changements de contexte volontaires, lorsque le parent se bloque dans \texttt{read}, et les involontaires, lorsqu'un enfant réveillé préempte le parent~; comme l'un peut remplacer l'autre, seul le total admet une borne inférieure nette.

Le coût direct tient à l'entrée dans le noyau, à l'exécution de l'ordonnanceur et au changement de registres. Entre processus, il change aussi la racine des tables de pages, ce qui peut faire perdre des traductions. Le coût le plus important vient en général après. L'autre tâche a chassé des caches les données de la première, ainsi que ses prédictions du prédicteur de branchement, si bien qu'elle tourne lentement un moment au retour. Ce coût dépend du travail effectué et n'apparaît pas dans un simple ping-pong. Les chiffres publiés varient beaucoup d'une machine à l'autre~: il faut donc les mesurer sur sa propre machine, et nous n'en donnons aucun ici.

Sur deux cœurs, chaque côté attendant activement sur une variable partagée sans se bloquer, aucun processus ne quitte son cœur, si bien qu'aucun changement de contexte n'a lieu. Le message traverse par un transfert de ligne de cache. Le prix à payer est deux cœurs à pleine charge, même au repos. C'est le compromis habituel du chemin critique d'un système de trading.

\iqwrong{«~Un appel système et un changement de contexte sont la même chose.~» Un appel système entre dans le noyau et revient en général au même thread, sans changement de tâche.}
\end{solution}

\begin{solution}{iq:iv3:operating-systems-concurrency-and-networks:15}
Le round robin de quantum 2 termine A à 13, B à 9 et C à 5, soit une moyenne de 9. Premier arrivé, premier servi termine à 8, 12 et 13, soit 11. Le plus court d'abord termine C à 1, B à 5 et A à 13, soit 6.33. Quand le quantum croît, le round robin devient premier arrivé, premier servi, ce que confirme le test.
\omcode{code/interviews-3/05-operating-systems-concurrency-and-networks/python/iv_round_robin.py}{1}{19}{Le round robin.\protect\hfil\protect\break}\label{lst:iv3:os:rr}

La trace pour le quantum 2 donne A jusqu'à 2, B jusqu'à 4, C jusqu'à 5, A jusqu'à 7, B jusqu'à 9, puis A jusqu'à 11 et 13. La dernière fin tombe toujours à 13, somme du travail, pour tout ordre qui ne laisse jamais le cœur inactif. Le test le vérifie et compare la fonction à un second modèle, tick par tick, sur 300 cas aléatoires.

Le plus court d'abord donne la plus faible moyenne des trois, et en fait la plus faible de tout ordre, car la durée de chaque tâche est comptée une fois pour chaque tâche qui la suit. Il pose deux problèmes. On ne connaît pas les durées à l'avance, et un flux continu de tâches courtes affame indéfiniment une tâche longue.

Le round robin n'a pas besoin de connaître les durées, et il borne le temps d'attente avant l'exécution. Une tâche courte se termine tôt, comme C ici. Un quantum très petit se rapproche d'un partage équitable du cœur, mais dans un système réel chaque changement de contexte coûte du temps, que le modèle ignore, et à partir d'un certain point le quantum n'est plus qu'un surcoût. Un ordonnanceur réel comme celui de Linux est plus complexe et suit le temps de processeur reçu par chaque tâche.

\iqwrong{«~Le round robin est le plus équitable, donc il est le meilleur pour le temps moyen d'achèvement.~» Sur des tâches de même durée, il fait moins bien que de les exécuter l'une après l'autre, car toutes se terminent vers la fin.}
\end{solution}

\begin{solution}{iq:iv3:operating-systems-concurrency-and-networks:16}
Le \Cref{lst:iv3:os:waitfor} suit chaque thread jusqu'à la ressource qu'il attend, puis jusqu'au détenteur de ce verrou, et ainsi de suite jusqu'à ce que la chaîne se répète ou s'arrête. Il faut que soient réunies~: l'exclusion mutuelle, la détention et l'attente, l'absence de préemption et l'attente circulaire. En pratique, on supprime en général l'attente circulaire, en prenant les verrous dans un ordre global unique.
\enlargethispage{-2\baselineskip}%
\omcode{code/interviews-3/05-operating-systems-concurrency-and-networks/python/iv_waitfor.py}{1}{17}{Suivi des attentes.\protect\hfil\protect\break}\label{lst:iv3:os:waitfor}

Aucune recherche générale n'est nécessaire, puisque chaque thread se bloque sur au plus un verrou, et que chaque verrou a un seul détenteur. Chaque thread n'a donc au plus qu'une arête sortante, et depuis n'importe quel point de départ il n'y a qu'un seul chemin à suivre. Le parcours s'arrête soit sur un thread en cours d'exécution, soit en entrant dans un cycle. \texttt{done} rend le travail total linéaire, puisqu'aucun thread n'est parcouru deux fois.

Les autres conditions peuvent elles aussi être levées. La détention et l'attente disparaît si l'on prend tous les verrous d'un coup, comme le fait \texttt{std::scoped\_lock}. L'absence de préemption tombe si l'on utilise un try-lock et que l'on recule en cas d'échec, au risque d'un livelock.

Le test compare la fonction à une recherche d'accessibilité sur 2000 états aléatoires, dont plus de 100 présentent un interblocage, et vérifie que chaque cycle renvoyé est bien un cycle du graphe.

Avec des verrous lecteurs-rédacteurs ou un sémaphore, un verrou peut avoir plusieurs détenteurs. Le graphe porte alors plusieurs arêtes par nœud et exige une vraie recherche de cycle~; et, pour un sémaphore, un cycle ne prouve plus un interblocage, car un autre détenteur peut encore libérer.

\iqwrong{«~Un thread bloqué depuis longtemps est en interblocage.~» Être bloqué derrière un détenteur lent n'est pas un interblocage. Il faut un cycle, et tous ceux qui s'y trouvent doivent attendre.}
\end{solution}

\begin{solution}{iq:iv3:operating-systems-concurrency-and-networks:17}
Le programme affiche \texttt{b\textquotesingle{}Saaa\textquotesingle{} b\textquotesingle{}Saaa\textquotesingle{} b\textquotesingle{}SPaa\textquotesingle{}} et \texttt{b\textquotesingle{}SaTa\textquotesingle{} b\textquotesingle{}SaTa\textquotesingle{} b\textquotesingle{}SPaa\textquotesingle{}}. Le manuel promet que les écritures d'un mappage partagé atteignent le fichier et les autres mappages, et que celles d'un mappage privé restent privées. Le fait que le mappage privé voie ou non les modifications ultérieures du fichier n'est pas précisé, et c'est de là que vient le \texttt{S}.

Le mappage privé a obtenu le \texttt{S}, puisque sa page était encore la page propre du fichier à cet instant, dans le cache de pages. Sa première écriture a copié la page. Par la suite, il était indépendant, et le \texttt{T} ultérieur n'apparaît donc pas. Une lecture par \texttt{pread} voit aussitôt l'écriture partagée, puisque les deux passent par un même cache de pages.

Un mappage convient à l'accès aléatoire sur un grand fichier, comme une base de ticks (tick store), puisque le noyau charge les données à la demande, les évince sous pression, et évite une copie dans un tampon utilisateur. Il permet aussi à plusieurs processus de partager une seule copie en lecture seule en mémoire. \texttt{read} est souvent préférable pour les parcours séquentiels. Il signale aussi une erreur d'E/S par une valeur de retour, alors qu'une lecture mappée échoue par un signal. Et sur le chemin critique, une lecture mappée peut se bloquer sur un défaut de page, sans aucun appel système en vue.

Une écriture partagée atteint le disque lorsque le noyau réécrit la page, un peu plus tard. Pour la forcer, il faut appeler \texttt{msync}, ou \texttt{fsync} sur le fichier.

\iqwrong{«~\texttt{mmap} est toujours plus rapide que \texttt{read}.~» Les défauts ont leur propre coût. Pour un parcours séquentiel, \texttt{read} avec lecture anticipée l'emporte souvent.}
\end{solution}

\begin{solution}{iq:iv3:operating-systems-concurrency-and-networks:18}
Un \texttt{write} réussi signifie que les octets ont atteint le cache de pages~; le disque peut donc contenir l'ancien fichier, le nouveau, ou un mélange quelconque des deux. Le \Cref{lst:iv3:os:atomicsave} écrit un nouveau fichier, le force sur le disque, le renomme par-dessus l'ancien nom, puis force le répertoire.
\omcode{code/interviews-3/05-operating-systems-concurrency-and-networks/python/iv_atomic_save.py}{1}{16}{Une sauvegarde atomique.\protect\hfil\protect\break}\label{lst:iv3:os:atomicsave}

Chaque appel a son rôle. \texttt{f.flush} transmet le tampon de Python au noyau. \texttt{fsync} sur le fichier force les données et leurs métadonnées vers le périphérique, et il doit venir avant le renommage~: sinon un plantage pourrait laisser le nouveau nom sur un fichier vide. \texttt{replace}, qui correspond à \texttt{rename(2)}, déplace le nom en une seule étape, si bien qu'un autre processus ne le trouve jamais absent. \texttt{fsync} sur le répertoire rend ce renommage durable, car un fsync sur un fichier ne sauvegarde pas forcément l'entrée de répertoire.

Le test enregistre l'ordre des appels, et simule un plantage avant le renommage, qui laisse l'ancien contenu intact. Aucun test ici ne coupe l'alimentation~; la durabilité repose donc sur les pages de manuel et sur le périphérique.

Ce qui reste sans protection~: un disque qui acquitte un flush qu'il n'a pas effectué, ce qui relève de la faute du matériel. Un plantage après le renommage et avant le flush du répertoire peut donner l'une ou l'autre version, chacune complète, et l'appelant doit pouvoir l'accepter. Les mises à jour entre deux sauvegardes sont perdues, car cela rend chaque sauvegarde atomique, et non chaque modification. C'est pourquoi un système qui ne doit rien perdre ajoute chaque modification à un journal, avec un fsync par validation, la sauvegarde n'étant qu'une optimisation par-dessus. Chaque fsync est lent, ce qui explique que de nombreux systèmes regroupent leurs vidages par lots.

\iqwrong{«~Le renommage est atomique, donc le fsync sur le fichier n'est pas nécessaire.~» Le renommage ordonne le changement de nom, pas les données. Le nom peut survivre à un plantage avec des données qui ne sont jamais parvenues sur le disque.}
\end{solution}

\begin{solution}{iq:iv3:operating-systems-concurrency-and-networks:19}
Ici, le programme affiche \texttt{1e-09} et \texttt{0}. Il faut utiliser \texttt{CLOCK\_MONOTONIC} pour chronométrer une requête, car il ne saute pas, et \texttt{CLOCK\_REALTIME} pour l'heure du jour, car seule celle-ci se rapporte à un calendrier et aux autres machines. Lire une horloge ne nécessite pas d'appel système, puisque le vDSO projette dans le processus le code et les données nécessaires à cette lecture.

\texttt{CLOCK\_REALTIME} peut sauter lorsque l'heure est réglée, si bien qu'une durée calculée à partir de lui peut être fausse, voire négative. \texttt{CLOCK\_MONOTONIC} ne saute pas, mais son rythme est légèrement ajusté par des services de type NTP. \texttt{CLOCK\_MONOTONIC\_RAW} échappe à ces ajustements, et il sert à comparer avec le matériel. Le temps monotone part d'une origine arbitraire, si bien qu'il ne dit rien à une autre machine.

Le vDSO existe parce que ces appels sont assez fréquents pour dominer le coût. Lire l'horloge coûte alors à peu près un appel de fonction, sans entrée dans le noyau. Nous avons vérifié que la projection est présente, dans \texttt{/proc/self/maps}, mais pas le comportement de chaque source d'horloge~; vérifiez donc ce que fait votre machine.

La résolution est le plus petit pas que l'horloge peut afficher, ici 1~ns. La justesse est la proximité d'une lecture avec le temps vrai, et un serveur peut s'en écarter bien davantage, selon la façon dont il est synchronisé. Deux horodatages venant de machines différentes peuvent donc être dans le désordre, quelle que soit leur résolution.

\iqwrong{«~Une résolution en nanosecondes signifie des horodatages exacts à la nanoseconde.~» Les derniers chiffres de l'horloge murale peuvent ne rien signifier.}
\end{solution}

\begin{solution}{iq:iv3:operating-systems-concurrency-and-networks:20}
Il affiche \texttt{connected and sent}, puis \texttt{b\textquotesingle{}early\textquotesingle{}}. C'est le noyau, et non votre processus, qui achève la poignée de main. \texttt{accept} prend simplement une connexion déjà établie dans une file, et les données arrivées en avance attendent dans le tampon de réception de cette connexion.

La poignée de main se déroule ainsi~: le client envoie SYN, le serveur répond par SYN et ACK, puis le client envoie ACK. \texttt{connect} rend la main après la deuxième de ces étapes. \texttt{listen} demande au noyau de faire cela pour le port. Le backlog fixe la taille de la file des connexions achevées qui attendent \texttt{accept}, limitée par un paramètre système (\texttt{somaxconn}).

Si la file est pleine, le manuel admet deux issues~: le client reçoit une erreur, ou sa requête est ignorée afin qu'une retransmission réussisse plus tard. TCP sous Linux prend la seconde par défaut, si bien qu'un client peut voir \texttt{connect} durer beaucoup plus longtemps. Un serveur serré, qui accepte lentement pendant une rafale de reconnexions, peut ainsi bloquer les clients. Le comportement diffère selon le système et dépend des paramètres.

Une connexion peut donc être établie et même contenir des données, dans un processus serveur qui n'en a jamais entendu parler. Si ce processus se termine ensuite, le client reçoit une remise à zéro. Cela compte pour les contrôles de santé~: une connexion TCP réussie montre que le noyau écoute, et elle ne montre pas que votre application est vivante.

\iqwrong{«~\texttt{accept} effectue la poignée de main.~» Elle a lieu avant, dans le noyau, et le processus peut être occupé ou bloqué pendant toute la durée.}
\end{solution}

\begin{solution}{iq:iv3:operating-systems-concurrency-and-networks:21}
Il affiche deux fois \texttt{(True, \textquotesingle{}EADDRINUSE\textquotesingle{})}, puis \texttt{(True, \textquotesingle{}bound\textquotesingle{})}. Le côté qui ferme en premier entre dans \texttt{TIME\_WAIT}, ici celui du serveur. La deuxième ligne échoue, car Linux exige \texttt{SO\_REUSEADDR} à la fois sur l'ancien et sur le nouveau socket~: le socket resté en \texttt{TIME\_WAIT} a été créé sans cette option.

C'est ce que cache d'ordinaire ce problème de redémarrage. L'ancien processus a disparu, sa connexion est en \texttt{TIME\_WAIT}, et le port paraît occupé un certain temps. Il faut donc placer \texttt{SO\_REUSEADDR} sur le socket d'écoute avant \texttt{bind}, dès la première version. Cela n'autorise pas deux écouteurs actifs sur une même adresse.

\texttt{TIME\_WAIT} remplit deux fonctions et dure deux fois la durée de vie maximale d'un segment. Le dernier ACK de la fermeture peut se perdre, auquel cas le pair renvoie son FIN, et ce côté doit encore être là pour répondre. Par ailleurs, d'anciens segments de cette connexion peuvent encore circuler dans le réseau, si bien que le même quadruplet ne doit pas être réutilisé tant qu'ils ne se sont pas éteints.

Cet état est donc correct, et il coûte peu~: un petit objet du noyau par connexion fermée. Il peut devenir un vrai problème pour un client qui ouvre un très grand nombre de connexions courtes vers un même serveur, et qui peut alors épuiser ses ports source. Le remède consiste en général à réutiliser les connexions. Réduire la durée de cet état ou le supprimer plus tôt retire une protection~; assurez-vous donc de l'avoir bien comprise d'abord.

\iqwrong{«~Fermer le socket libère le port immédiatement.~» Seulement si ce côté n'a pas été celui qui a fermé en premier, ou s'il est passé par \texttt{TIME\_WAIT}.}
\iqfollow{Pourquoi le côté client n'a-t-il pas ce problème~?}{En général, il l'a aussi, mais son port change à chaque connexion, si bien que personne ne le remarque. Ici, le serveur a fermé en premier, donc c'est lui qui paie.}
\end{solution}

\begin{solution}{iq:iv3:operating-systems-concurrency-and-networks:22}
Le produit vaut $10^{10} \times 10^{-3} / 8 = 1.25$~Mo. Avec une fenêtre de 65\,535~octets, le débit vaut $65\,535 \times 8 / 10^{-3} = 524$~Mbit/s, soit environ 5~\% de la liaison. Le facteur d'échelle de fenêtre multiplie le champ de 16~bits par une puissance de 2, convenue lors de la poignée de main, si bien que la fenêtre peut atteindre 1~Gio.
\omcode{code/interviews-3/05-operating-systems-concurrency-and-networks/python/iv_window.py}{1}{15}{Fenêtre, temps aller-retour et débit.\protect\hfil\protect\break}\label{lst:iv3:os:window}

Un émetteur ne peut avoir en vol qu'une fenêtre de données non acquittées au plus, et les acquittements prennent un aller-retour. Une fenêtre part donc à chaque aller-retour, quelle que soit la vitesse de la liaison. Pour remplir la liaison, la fenêtre doit atteindre le produit débit-délai. Le test compare la formule à un décompte aller-retour par aller-retour sur 200 cas aléatoires, et il vérifie les deux extrêmes~: une petite fenêtre limite le débit, et une grande fenêtre ne laisse plus que la liaison comme limite.

Si l'on double le temps aller-retour avec une fenêtre fixe, le débit est divisé par deux. Le cas extrême est une liaison intercontinentale, où une petite fenêtre par défaut donne un faible débit sur un chemin pourtant très rapide. C'est pourquoi les réglages des tampons comptent davantage à longue distance.

Les deux mécanismes répondent à des questions différentes. Le contrôle de flux protège le récepteur~: il annonce la taille de son tampon libre, et c'est la fenêtre de l'en-tête. Le contrôle de congestion protège le réseau~: l'émetteur tient sa propre fenêtre de congestion, la fait croître, et la réduit à l'apparition de signes de perte ou de délai. L'émetteur utilise la plus petite des deux.

La réponse ignore les pertes. Un segment perdu conduit un émetteur ordinaire à réduire son débit, si bien que sur un long chemin, même un faible taux de perte peut faire chuter le débit bien en dessous de cette limite.

\iqwrong{«~Une liaison plus rapide corrige un transfert lent.~» Si la fenêtre limite le transfert, une liaison plus rapide ne change rien.}
\end{solution}

\begin{solution}{iq:iv3:operating-systems-concurrency-and-networks:23}
\Cref{lst:iv3:os:framing} garde un tampon, et ne retire que des messages complets. \texttt{recv} rend ce qui est arrivé, qui peut être une partie d'un message, plusieurs messages, ou un mélange des deux. En boucle locale, une petite écriture arrive en général entière, si bien que la mauvaise version passe~; sur un vrai réseau, avec segmentation et rafales, elle ne passe pas.
\enlargethispage{-3\baselineskip}%
\omcode{code/interviews-3/05-operating-systems-concurrency-and-networks/python/iv_framing.py}{1}{26}{Découpage d'un flux d'octets en messages.\protect\hfil\protect\break}\label{lst:iv3:os:framing}

TCP ne sait pas où commencent vos messages. Le champ de longueur vient en premier, si bien que le lecteur sait toujours combien d'octets il lui faut, et la boucle traite plusieurs messages en une seule lecture. Si seule une partie d'un message est là, elle reste dans le tampon jusqu'à l'appel suivant.

La limite de longueur compte. Sans elle, une longueur corrompue fait attendre le lecteur jusqu'à 65\,535~octets qui ne viendront peut-être jamais, puis livrer des données inutilisables comme un message, toutes les trames suivantes étant décalées. Avec un champ de longueur plus large, la même faute pourrait retenir des gigaoctets. Un corps vide est un message légal, et le test le couvre, car un \texttt{if not body} imprudent le perdrait. Le test fait aussi passer 300 messages par des découpages aléatoires, puis sur une vraie connexion TCP lue par morceaux de 97~octets.

\texttt{del self.buf[:n]} est bon marché dans CPython, car un \texttt{bytearray} ne fait que déplacer son début logique, et le test voit l'adresse du premier octet avancer de $n$. Dans un langage où cela décale les octets, garder une position de lecture, et compacter de temps à autre.

Du côté de l'envoi, \texttt{send} peut n'accepter qu'une partie des données~; utiliser donc \texttt{sendall} sur un socket bloquant, ou garder le reste et réessayer sur un socket non bloquant. Deux threads ne doivent pas écrire sur un même socket sans verrou, sans quoi leurs messages peuvent s'entremêler.

\iqwrong{«~TCP est fiable, donc un \texttt{send} donne un \texttt{recv}.~» La fiabilité concerne les octets. Rien ne préserve les frontières.}
\end{solution}

\begin{solution}{iq:iv3:operating-systems-concurrency-and-networks:24}
Il affiche \texttt{b\textquotesingle{}abcd\textquotesingle{} b\textquotesingle{}ij\textquotesingle{}}, puis \texttt{65507}, puis \texttt{Message too long}. Les octets \texttt{efgh} ont été perdus dès que la lecture de 4~octets a pris le premier datagramme. Chaque \texttt{recv} consomme un datagramme entier, quelle que soit la taille du tampon. 65\,507 correspond à la limite de 65\,535 du champ de longueur IP, moins 20~octets d'en-tête IP et 8~d'en-tête UDP.

Rien ne signale la troncature par défaut. \texttt{recvmsg} place un indicateur \texttt{MSG\_TRUNC} dans ses résultats, et, avec cet indicateur passé à \texttt{recv}, l'appel rend la longueur réelle. Un gestionnaire de flux doit donc utiliser un tampon plus grand que le plus gros datagramme qu'il peut recevoir, ou vérifier cet indicateur.

Un flux réel garde ses datagrammes sous le MTU du chemin, typiquement en dessous d'environ 1500~octets sur Ethernet. Les plus grands sont découpés en fragments IP, et si un fragment est perdu, le datagramme entier l'est aussi, si bien que le taux de perte se multiplie. Les fragments ajoutent aussi du travail de réassemblage. Les plateformes précisent la taille de leurs paquets, lire donc la spécification et dimensionner le tampon en conséquence.

Sur une interface de boucle locale, le MTU est grand, ce qui explique que le maximum fonctionne ici sans problème visible. Cela ne dit rien d'un vrai réseau.

\iqwrong{«~Le reste du datagramme reste pour le \texttt{recv} suivant.~» C'est le comportement d'un flux TCP. En UDP, la lecture suivante obtient le datagramme suivant.}
\end{solution}

\begin{solution}{iq:iv3:operating-systems-concurrency-and-networks:25}
Il affiche \texttt{8192}, \texttt{sent 1000, no error}, \texttt{True} et \texttt{[True]}. Le noyau double la taille demandée pour tenir compte de sa propre comptabilité, si bien que 4096 se relit comme 8192. Le tampon du socket de réception a débordé, le noyau a jeté le reste à cet endroit, et la dernière colonne de \texttt{/proc/net/udp} compte ces pertes.

L'émetteur ne voit rien, car UDP n'a pas de retour. Un \texttt{sendto} réussi signifie seulement que le paquet a été accepté pour l'envoi. Rien n'avertit non plus le programme récepteur, et c'est pourquoi il faut surveiller les compteurs. Le nombre de datagrammes reçus est inférieur à 1000 sur toute exécution où le tampon est aussi petit et où personne ne lit pendant l'envoi.

Pour un gestionnaire de flux qui perd des données lors des rafales, l'ordre des actions est le suivant. Il faut d'abord mesurer où a lieu la perte~: tampons du socket, anneau de la carte réseau, ou réseau. Les numéros de séquence montrent qu'une donnée a été perdue, et les compteurs indiquent où. Il faut ensuite regarder le lecteur. Sur son thread, il ne doit faire que lire et transmettre, les travaux lents étant déplacés ailleurs. Donner-lui un cœur à lui seul, et garder les interruptions et le lecteur sur le bon nœud NUMA. Augmenter ensuite \texttt{SO\_RCVBUF}, en relevant si besoin la limite système \texttt{rmem\_max}. Un tampon plus grand ne fait qu'absorber une rafale, et si le lecteur est en moyenne plus lent, il perd encore des données, simplement plus tard. Pour l'étape finale, envisager le contournement du noyau.

\iqwrong{«~Rendez le tampon immense et les pertes disparaissent.~» Cela aide contre les rafales. Cela ajoute du délai, et ne change rien pour un lecteur trop lent dans l'ensemble.}
\end{solution}

\begin{solution}{iq:iv3:operating-systems-concurrency-and-networks:26}
Un groupe est une adresse de 224.0.0.0/4 que n'importe quelle machine peut rejoindre. Un émetteur, qui n'a pas besoin de rejoindre le groupe, envoie une seule copie, et le réseau la recopie vers tous les membres. Une machine rejoint le groupe avec \texttt{IP\_ADD\_MEMBERSHIP}. Avec \texttt{join}, le listing affiche \texttt{b\textquotesingle{}tick\textquotesingle{}}, et sinon \texttt{nothing}. 239.1.2.3 correspond à \texttt{01:00:5e:01:02:03}, et 31 autres groupes partagent cette adresse.
\omcode{code/interviews-3/05-operating-systems-concurrency-and-networks/python/iv_mcast_mac.py}{1}{7}{Adresse de groupe vers adresse Ethernet.\protect\hfil\protect\break}\label{lst:iv3:os:mcastmac}

L'adresse Ethernet ne garde que les 23~bits de poids faible, et un groupe compte 28~bits significatifs, donc 5~bits sont perdus, et $2^5 = 32$ groupes correspondent à chaque adresse. 224.1.2.3 et 239.129.2.3 en font partie. La carte réseau peut donc transmettre au système des paquets d'un groupe que vous n'avez pas rejoint, et le noyau les filtre. Les flux très chargés doivent choisir des adresses de groupe qui n'entrent pas en collision.

Rejoindre un groupe a aussi un versant réseau. La machine envoie un rapport IGMP, et les commutateurs dotés d'IGMP snooping s'en servent pour n'envoyer le groupe qu'aux ports qui l'ont demandé. L'appartenance est propre à chaque interface, si bien que désigner la mauvaise est une cause classique de ne rien recevoir sur une machine dotée de plusieurs réseaux.

\texttt{IP\_MULTICAST\_TTL} vaut 1 par défaut, si bien que le paquet ne traverse pas de routeur, et \texttt{IP\_MULTICAST\_LOOP} est activé par défaut, si bien que la machine de l'émetteur reçoit elle-même le paquet. Nous avons lu les deux valeurs sur un socket neuf. Un test sur une seule machine peut donc fonctionner, alors que deux machines sur des sous-réseaux différents ne reçoivent rien, à moins que le TTL ne soit relevé et que les routeurs ne soient configurés pour le multicast.

L'exécution ici a lieu en boucle locale, sans commutateur ni IGMP sur un câble, si bien que ce que fait le réseau repose sur les sources inscrites au registre.

\iqwrong{«~Il suffit de se lier au port du groupe.~» Sans rejoindre le groupe, le noyau n'a aucune raison d'accepter ce groupe, comme le montre la seconde exécution.}
\end{solution}

\begin{solution}{iq:iv3:operating-systems-concurrency-and-networks:27}
Il affiche \texttt{(\textquotesingle{}EINPROGRESS\textquotesingle{}, 0)} et \texttt{(\textquotesingle{}EINPROGRESS\textquotesingle{}, \textquotesingle{}ECONNREFUSED\textquotesingle{})}. Un \texttt{connect} non bloquant lance la poignée de main et rend la main aussitôt, si bien que son résultat signifie seulement «~en cours~», même lorsque la tentative va échouer. Le résultat réel est livré plus tard, dans \texttt{SO\_ERROR}.

Un socket accessible en écriture n'est donc pas connecté. \texttt{select} marque le socket accessible en écriture lorsque la tentative est terminée, dans l'un ou l'autre cas. Si vous le prenez pour un succès et commencez à envoyer, vous obtenez une erreur déroutante à la première écriture, loin de la cause. Lisez \texttt{SO\_ERROR} au réveil~: 0 signifie connecté, et toute autre valeur donne la raison de l'échec. La lecture de l'option la remet aussi à zéro.

Le premier cas du listing est un serveur qui existe, et le second un port fermé. Tous deux renvoient ici \texttt{EINPROGRESS}, même en boucle locale, et c'est ce qui rend la vérification nécessaire. Le manuel admet d'autres résultats, il faut donc aussi traiter le cas d'un retour à 0, qui signifie une connexion immédiate.

Un connect bloquant peut attendre longtemps sur un hôte injoignable, avec plusieurs retransmissions du SYN. Un processus qui doit rester réactif se connecte donc sans bloquer, et fixe lui-même sa limite de temps à l'aide de \texttt{select} ou d'\texttt{epoll}.

\iqwrong{«~\texttt{select} signale la socket accessible en écriture, donc nous sommes connectés.~» Cela signifie que la tentative est terminée.}
\iqfollow{Que devez-vous faire en cas d'échec~?}{On ferme cette socket et on en crée une nouvelle, car un connect échoué laisse la socket dans un état non spécifié sur certains systèmes. On réessaie avec un délai croissant.}
\end{solution}

\begin{solution}{iq:iv3:operating-systems-concurrency-and-networks:28}
Sans argument, le programme C++ est tué par \texttt{SIGPIPE} et n'affiche rien. Avec un argument, \texttt{send} utilise \texttt{MSG\_NOSIGNAL}, si bien qu'il renvoie $-1$ et que le programme affiche \texttt{-1 Broken pipe}. Dans le listing Python, le premier send renvoie 3, et le second lève \texttt{BrokenPipeError}, le nom que Python donne à \texttt{EPIPE}.

Un processus qui écrit dans une connexion sans lecteur reçoit un signal, et par défaut ce signal met fin au processus. Un serveur peut donc mourir parce qu'un client s'en est allé, à moins de prendre des précautions. Il faut mettre \texttt{SIGPIPE} en mode ignoré au démarrage, ou passer \texttt{MSG\_NOSIGNAL} à chaque send, puis traiter la valeur de retour d'erreur. Python le met lui-même en ignoré, ce qui explique l'exception qu'on observe ici.

Le premier send réussit, car à ce stade ce côté sait seulement que le pair a envoyé un FIN. Un FIN signifie que le pair n'a plus rien à émettre, mais il peut encore lire, si bien que TCP accepte les données et les envoie. Le noyau du pair, constatant que la socket a été fermée et que personne ne peut lire ces données, répond par un RST. Une fois celui-ci arrivé, l'écriture suivante échoue. Le listing attend cette réponse avec un \texttt{poll} qui ne demande aucun événement, car seule une erreur ou une déconnexion le réveille alors. Un \texttt{select} en lecture n'attendrait pas, puisque le FIN rend déjà la socket lisible. Sur un réseau, sans attente, le second send pourrait passer en premier et réussir aussi.

Un \texttt{send} réussi ne signifiait donc jamais que le pair avait lu les données. Des octets encore présents dans le tampon local peuvent être perdus, et si l'on doit le savoir, l'application a besoin de son propre accusé de réception. Pour la saisie d'ordres, c'est pourquoi une session dispose de numéros de séquence et de réponses.

\iqwrong{«~Une connexion fermée fait échouer l'écriture suivante immédiatement.~» Seulement après l'arrivée du RST, ce qui demande un aller-retour.}
\end{solution}

\begin{solution}{iq:iv3:operating-systems-concurrency-and-networks:29}
La tâche de haute priorité se termine à 26. Il s'agit d'une inversion de priorité~: la tâche intermédiaire, qui n'a besoin d'aucun verrou, retarde une tâche plus prioritaire de la totalité de ses 20 ticks. Avec l'héritage de priorité, le détenteur du verrou s'exécute à la priorité de celui qui l'attend, et la tâche de haute priorité se termine à 6.
\omcode{code/interviews-3/05-operating-systems-concurrency-and-networks/python/iv_inversion.py}{1}{34}{Un processeur et un verrou.\protect\hfil\protect\break}\label{lst:iv3:os:inversion}

La trace sans héritage~: la tâche basse s'exécute de 0 à 1 et détient le verrou. La tâche haute arrive et se bloque aussitôt sur le verrou, si bien que la tâche basse poursuit jusqu'à 2, moment où la tâche intermédiaire arrive et prend le cœur jusqu'à 22. Ce n'est qu'alors que la tâche basse peut se terminer, à 24. La tâche haute prend ensuite 2 ticks. Avec héritage, la tâche basse s'exécute jusqu'à 4 à haute priorité, et la tâche haute se termine à 6, comme si la tâche intermédiaire n'existait pas.

Le test montre cette forme. Sans héritage, la tâche haute se termine à $6 + m$ pour une tâche intermédiaire de $m$ ticks, si bien que le retard croît sans limite. Avec l'héritage, la durée vaut toujours 6, bornée par ce qui reste de la section critique. Cela est vérifié sur 200 cas aléatoires. La date de fin de la dernière tâche ne change pas avec l'héritage, car le travail total est le même~: seul change l'identité de celui qui attend.

Les autres remèdes~: garder les sections critiques courtes, éviter de partager des verrous entre priorités, ou faire exécuter tous les détenteurs de verrous à une priorité plafond relevée. Linux dispose de \texttt{PTHREAD\_PRIO\_INHERIT} pour les mutex. Un système de trading qui utilise des priorités temps réel sur des cœurs épinglés court le même risque~: un thread de basse priorité qui détient un verrou et a perdu son cœur peut bloquer le chemin critique (hot path). Un thread \texttt{SCHED\_FIFO} qui tourne en attente active sur un verrou détenu par une priorité temps réel inférieure sur son propre cœur attend indéfiniment, car rien ne permet au détenteur de s'exécuter. Si le détenteur est une tâche normale, il ne s'exécute que dans la tranche que le noyau réserve aux tâches temps réel, soit 5~\% par défaut.

\iqwrong{«~Donner à la tâche importante la priorité la plus haute, et elle s'exécutera toujours en premier.~» La priorité ordonne les tâches exécutables. Une tâche bloquée sur un verrou n'est pas exécutable.}
\end{solution}

\begin{solution}{iq:iv3:operating-systems-concurrency-and-networks:30}
Le \Cref{lst:iv3:os:rwlock} en donne une réponse. \texttt{may\_read} vaut faux dès qu'un écrivain attend. Sans cette préférence, un flux continu de lecteurs qui se recouvrent pourrait repousser un écrivain indéfiniment. Le risque se déplace simplement~: un flux continu d'écrivains peut désormais provoquer la famine des lecteurs.
\omcode{code/interviews-3/05-operating-systems-concurrency-and-networks/cpp/rwlock.hpp}{5}{42}{Un verrou lecteurs-écrivains qui privilégie les écrivains.\protect\hfil\protect\break}\label{lst:iv3:os:rwlock}

Le dernier lecteur réveille un écrivain, car le prédicat de l'écrivain ne peut pas être vrai avant cela. Un écrivain qui sort réveille l'écrivain suivant s'il y en a un, et sinon tous les lecteurs, puisque plusieurs peuvent entrer à la fois. Tous les changements d'état se font sous le mutex, si bien qu'aucun réveil ne peut être perdu.

Le test de préférence n'a besoin d'aucune hypothèse de temps dans son assertion. Un lecteur détient le verrou. Un thread écrivain appelle \texttt{lock} et se bloque, et le pilote ne dort que pour lui laisser le temps d'arriver. Alors \texttt{try\_lock\_shared} doit renvoyer faux. Si l'écrivain n'était pas encore arrivé, l'appel renverrait vrai et le test échouerait, il ne réussirait jamais à tort. Un jumeau sans le contrôle de l'écrivain affiche 1 sur cette ligne, si bien que le test peut prendre en défaut le modèle. Sous ThreadSanitizer, les écrivains et les lecteurs s'exécutent avec un invariant ($a = b$ sur des variables ordinaires), et aucun rapport de course aux données n'apparaît.

Il se peut qu'il ne soit pas plus rapide qu'un mutex, car un lecteur doit tout de même écrire l'état du verrou, en prenant ici un mutex interne. Ainsi, sur de nombreux cœurs, tous les lecteurs se disputent une même ligne de cache, et pour une section critique très courte, ce trafic constitue le principal coût. Un verrou en lecture ne devient rentable que lorsque les lecteurs le détiennent assez longtemps pour qu'un véritable recouvrement ait lieu. Nous n'avons rien mesuré ici. D'autres options, pour des lectures très brèves et des écritures rares, sont le verrou à séquence (seqlock), ou le remplacement d'un pointeur vers une copie immuable.

\iqwrong{«~Les lecteurs ne se gênent pas, donc un verrou en lecture passe à l'échelle.~» Leur travail n'entre pas en conflit. Ce sont leurs mises à jour du verrou qui le font.}
\end{solution}

\begin{solution}{iq:iv3:operating-systems-concurrency-and-networks:31}
Le \Cref{lst:iv3:os:futex} garde 0 pour libre, 1 pour verrouillé, et 2 pour verrouillé avec des attentes possibles. Verrouiller et déverrouiller sans contention ne sont que des opérations atomiques, et le million de tours du pilote ne fait aucun appel à \texttt{futex}. Trois états sont nécessaires pour que \texttt{unlock} puisse éviter l'appel système lorsque personne ne dort.
\omcode{code/interviews-3/05-operating-systems-concurrency-and-networks/cpp/futex_mutex.hpp}{7}{34}{Un mutex sur un futex.\protect\hfil\protect\break}\label{lst:iv3:os:futex}

Avec deux états, \texttt{unlock} ne peut pas savoir si quelqu'un dort, si bien qu'il doit soit faire un appel système à chaque fois, soit oublier des threads endormis. L'état 2 indique que quelqu'un peut dormir. Un thread qui n'obtient pas le verrou remplace la valeur par 2, puis s'endort. Si l'ancienne valeur était 0, son échange a en fait pris le verrou.

Un thread réveillé ne sait pas si d'autres dorment encore, il doit donc poser 2 et non 1. Avec 1, il libérerait ensuite le verrou sans réveil, et tout autre thread endormi resterait bloqué à jamais. Notre test construit ce jumeau, et le pilote, avec deux threads en attente derrière un détenteur, se bloque et atteint sa limite de temps. L'état 2 provoque parfois un réveil superflu, ce qui est inoffensif.

\texttt{FUTEX\_WAIT} prend la valeur attendue à cause d'une condition de course. Entre la décision d'un thread de dormir et son entrée dans le noyau, le détenteur peut déverrouiller et appeler le réveil, et comme personne ne dort encore, ce réveil est perdu. Le noyau vérifie donc le mot et bloque de manière atomique, uniquement si le mot vaut encore 2. Sinon, l'appel revient aussitôt. Cette vérification comble l'écart.

Les ordres mémoire sont acquire à la prise du verrou et release à sa libération, si bien que les accès à la section critique ne peuvent pas en sortir. Le compteur donne exactement 400\,000, et ThreadSanitizer ne signale aucune course aux données. Un vrai mutex ajoute une attente active avant de dormir, ainsi qu'une gestion très soignée, donc il faut utiliser \texttt{std::mutex} en production.

\iqwrong{«~Un mutex entre toujours dans le noyau.~» Seulement lorsqu'il doit dormir ou réveiller quelqu'un.}
\end{solution}

\begin{solution}{iq:iv3:operating-systems-concurrency-and-networks:32}
Environ 40.2~ms, soit 400 fois l'aller-retour. L'en-tête part aussitôt, mais Nagle retient le corps jusqu'à ce que l'en-tête soit acquitté, et le serveur retarde cet acquittement, car il n'a aucune réponse à envoyer avant d'avoir le corps. Chacun attend l'autre jusqu'à ce que la minuterie de 40~ms y mette fin.
\omcode{code/interviews-3/05-operating-systems-concurrency-and-networks/python/iv_nagle.py}{1}{38}{Écriture, écriture, lecture, sous forme d'événements.\protect\hfil\protect\break}\label{lst:iv3:os:nagle}

Chaque règle est raisonnable prise seule. La règle de Nagle n'envoie de petites données que lorsqu'aucune n'est en attente d'acquittement, ce qui a évité aux réseaux des déluges de minuscules paquets. L'acquittement différé attend un peu dans l'espoir de transporter l'ACK dans une réponse. La norme autorise ce délai jusqu'à 0.5~s, et l'énoncé prend 40~ms. Ensemble, dans le schéma écriture, écriture, lecture, elles produisent un interblocage jusqu'au déclenchement d'une minuterie. La chronologie est la suivante~: l'en-tête arrive à 0.05, la minuterie se termine à 40.05, l'ACK arrive à 40.1, le corps arrive à 40.15, la réponse arrive à 40.2.

Trois correctifs. La meilleure solution consiste à n'effectuer qu'une seule écriture~: placer l'en-tête et le corps dans un même tampon, ou utiliser \texttt{writev}, ce qui économise aussi un paquet et un appel système. Le deuxième est \texttt{TCP\_NODELAY}, qui envoie tout sans délai, soit 0.1~ms dans le modèle, et qui est le réglage habituel pour les connexions de trading, où l'application décide ce qu'est un message. Le troisième est \texttt{TCP\_QUICKACK} côté récepteur, qui donne 0.2~ms dans le modèle, mais il n'est pas permanent et doit être réappliqué.

Une seule écriture ne bloque jamais, comme le montre le test, car rien n'est en attente lorsqu'elle part. À mesure que la minuterie approche de 0, le blocage se réduit à un aller-retour supplémentaire. \texttt{TCP\_NODELAY} vaut 0 par défaut sur une nouvelle socket, ce que nous avons relu. Il s'agit d'un modèle, aucune capture de paquets n'a été réalisée, si bien que la valeur réelle de la minuterie reste à déterminer sur la machine.

\iqwrong{«~Désactiver Nagle rend TCP peu fiable, ou lui fait cesser les retransmissions.~» Cela arrête seulement de retenir les petits segments.}
\end{solution}

\begin{solution}{iq:iv3:operating-systems-concurrency-and-networks:33}
La carte extrait la trame du câble et l'écrit par DMA dans un tampon circulaire de réception en mémoire. Elle lève une interruption, à moins que le pilote ne scrute déjà. NAPI scrute dans le contexte d'interruption logicielle, remonte le paquet à travers IP et UDP, trouve la socket, et ajoute le paquet à sa file de réception. L'ordonnanceur réveille le thread, et \texttt{recv} copie les octets dans le tampon.

Étape par étape, voici où le paquet attend, est abandonné ou est copié.

La carte~: elle calcule un haché de l'en-tête pour choisir l'une de ses files (répartition côté réception (RSS)), si bien qu'un flux va vers une file et une interruption fixes. Si le tampon circulaire est plein, par exemple lorsque le noyau prend du retard, la carte abandonne le paquet. \texttt{ethtool -S} affiche ces compteurs. La modération des interruptions peut retenir l'interruption un court instant, afin de regrouper les paquets.

Le softirq~: les paquets sont traités sur un cœur, éventuellement choisi par l'orientation logicielle de RPS. Si cette file par cœur dépasse sa limite (\texttt{netdev\_max\_backlog}), le paquet est abandonné, et la deuxième colonne de \texttt{/proc/net/softnet\_stat} sur ce noyau en conserve le compte. Si le cœur est occupé par d'autres travaux, le paquet attend ici.

La socket~: UDP trouve la socket par adresse et port, et l'appartenance à un groupe multicast est vérifiée. Si le tampon de réception de la socket est plein, le paquet est abandonné, et \texttt{/proc/net/udp} et \texttt{nstat} le montrent.

Le réveil~: le thread doit être ordonnancé, sur un cœur qui peut être endormi ou occupé. Un thread qui tourne en attente active dans un \texttt{recv} non bloquant supprime cette attente. Vient ensuite la copie vers l'espace utilisateur, une copie de plus, et le coût de l'appel système.

Le contournement du noyau (kernel bypass) projette directement les tampons circulaires de la carte dans le processus, et le thread les scrute. Ainsi disparaissent l'interruption, le softirq, la pile réseau, la file de la socket, le réveil et la copie. Il reste le DMA, la scrutation, et un tampon circulaire plein si l'on est lent. On renonce aux outils du noyau, et il faut prendre en charge soi-même le traitement du protocole.

Cette réponse est une description, aucun paquet n'a été tracé pour elle, si bien que les étapes reposent sur les documents du noyau consignés dans le registre.

\iqwrong{«~Le paquet va de la carte directement à ma socket.~» Il y a plusieurs files intermédiaires, et chacune peut abandonner le paquet.}
\end{solution}

\begin{solution}{iq:iv3:operating-systems-concurrency-and-networks:34}
Il affiche \texttt{[1, 1, 1]} et \texttt{[1, 0, 1]}. Le mode déclenché par niveau signale la socket tant que des données restent non lues, si bien que les six octets restants donnent un événement au second poll. Le mode déclenché par front ne signale que les changements~: rien de nouveau n'est arrivé, donc le second poll est vide, et le troisième se déclenche puisque de nouvelles données sont arrivées.

Ainsi, en mode déclenché par front, si l'on ne lit qu'une partie des données et qu'on attend, on peut attendre indéfiniment avec des octets dans le tampon, si rien de nouveau n'arrive. La règle consiste à lire jusqu'à \texttt{EAGAIN}. Cela suppose une socket non bloquante, car une lecture bloquante sur une socket vide bloquerait toute la boucle. Il en va de même pour les écritures.

Le mode déclenché par front évite les réveils répétés pour une socket que l'on sait ne pas pouvoir servir encore, et il permet à plusieurs threads de partager un ensemble de sockets sans qu'ils se réveillent tous sur un même événement. Il est plus facile de s'y tromper~: le mode déclenché par niveau est la valeur par défaut la plus sûre.

\texttt{select} et \texttt{poll} reçoivent tout l'ensemble des descripteurs à chaque appel, et le noyau le parcourt, si bien que chaque appel coûte en proportion du nombre de descripteurs surveillés, quel que soit le nombre de descripteurs actifs. \texttt{epoll} enregistre une fois et tient une liste des descripteurs prêts, si bien que chaque appel coûte en proportion du nombre de descripteurs prêts. Pour dix mille connexions inactives, c'est une économie importante.

Il n'y a alors aucun gain réel lorsque presque tous les descripteurs sont actifs à chaque appel, ou lorsqu'il n'y en a qu'une poignée. Un chemin critique qui surveille quelques sockets peut ne jamais attendre, mais tourner en attente active avec des lectures non bloquantes, et éviter l'appel système que représenterait \texttt{epoll\_wait}.

\iqwrong{«~Le mode déclenché par front est plus rapide, donc on l'utilise toujours.~» Il évite les événements répétés, mais coûte au moins une lecture supplémentaire par rafale, pour constater \texttt{EAGAIN}. Les bogues dans ce mode bloquent le programme en silence.}
\end{solution}

\begin{solution}{iq:iv3:operating-systems-concurrency-and-networks:35}
Il affiche \texttt{parent threads 2}, \texttt{child threads 1, try\_lock 0} et \texttt{parent try\_lock 1}. Seul le thread appelant existe dans le fils. La mémoire est copiée telle quelle, y compris le mutex dans son état verrouillé, et le thread qui le déverrouillerait n'existe pas dans le fils. Le verrou du fils reste donc bloqué à jamais.

Ici \texttt{try\_lock} renvoie 0. Un simple \texttt{lock} dans le fils ne reviendrait jamais, et le programme se bloquerait sans aucune erreur pour le signaler. C'est la défaillance décrite dans l'énoncé~: un fils qui appelle une bibliothèque de journalisation dont le verrou interne était détenu au moment du fork se bloque au hasard, selon ce que faisaient les autres threads à ce moment-là. Il peut fonctionner pendant des mois, puis se bloquer.

Il en va de même pour les données protégées par le verrou. L'autre thread pouvait être à mi-chemin d'une mise à jour, si bien que la structure dans le fils est incohérente, et aucun thread n'est là pour la terminer. Certaines bibliothèques protègent leur état à l'aide de gestionnaires \texttt{pthread\_atfork}, mais on ne peut pas compter sur toutes les bibliothèques.

Ainsi, tant qu'il n'appelle pas \texttt{exec}, un fils d'un processus multithread ne peut sans risque appeler que des fonctions async-signal-safe. En pratique, on ne fait presque rien et on appelle aussitôt \texttt{exec}, ou on utilise \texttt{posix\_spawn}. Si l'on fait un fork pour prendre un instantané, il faut le faire avant de lancer le moindre thread, ou copier les données par un autre moyen.

\iqwrong{«~Le fils reçoit tous les threads, puisqu'ils font partie du processus.~» Il n'en reçoit qu'un. Tout ce que faisaient les autres s'arrête simplement, dans la copie du fils.}
\end{solution}

\begin{solution}{iq:iv3:operating-systems-concurrency-and-networks:36}
Sans argument, il affiche \texttt{EOWNERDEAD} et \texttt{0}~: le verrou est accordé, avec l'information que son propriétaire est mort. Avec un argument, un mutex ordinaire reste verrouillé, et le verrou temporisé renvoie \texttt{ETIMEDOUT}. Sans limite de temps, le survivant resterait bloqué à jamais, sans rien dans les journaux.

Le noyau peut le faire, car le propriétaire d'un mutex robuste tient une liste des verrous robustes qu'il détient. Quand un thread meurt, le noyau parcourt cette liste, marque chaque verrou d'un drapeau et réveille un thread en attente. Un mutex ordinaire n'est qu'un mot en mémoire, et personne ne sait qui le détient. Ainsi, le plantage d'un processus laisse ce mot verrouillé, et tout le segment partagé reste bloqué.

Sur \texttt{EOWNERDEAD}, le nouveau propriétaire détient le verrou, et les données qu'il protège peuvent être à moitié mises à jour, puisque le propriétaire est mort à une instruction quelconque. Seule l'application sait à quoi ressemble un état correct. Le nouveau propriétaire doit donc vérifier ou réparer les données, puis appeler \texttt{pthread\_mutex\_consistent}. S'il déverrouille sans cet appel, le mutex devient inutilisable de façon définitive, et tous les appels de verrouillage suivants renvoient \texttt{ENOTRECOVERABLE}. La bibliothèque ne peut pas réparer les données à la place de l'application, puisqu'elle ignore ce que protège le verrou.

La robustesse a donc un prix~: chaque mise à jour doit laisser les données réparables, par exemple en écrivant d'abord dans un journal, ou en marquant un enregistrement comme modifié pendant la mise à jour. En pratique, c'est la conception des données qui est difficile.

L'autre solution consiste à éviter tout verrou entre processus. Un tampon circulaire sans verrou (lock-free) permet à un producteur de mourir sans détenir de verrou, et un processus distinct peut surveiller cette mort avec \texttt{waitpid} ou un heartbeat.

\iqwrong{«~Le noyau libère les verrous d'un processus mort.~» Il libère sa mémoire et ses fichiers. Un mutex ordinaire en mémoire partagée n'appartient à personne.}
\end{solution}

\begin{solution}{iq:iv3:operating-systems-concurrency-and-networks:37}
Le \Cref{lst:iv3:os:arb} garde le prochain numéro attendu et un ensemble des numéros suivants. Un doublon ou un numéro ancien est rejeté, et tout ce qui devient contigu sort. Lorsque plus de \texttt{hold} numéros attendent, la plage manquante est signalée et ignorée. Attendre indéfiniment figerait le carnet d'ordres~: si les deux lignes ont perdu un paquet, il n'arrivera jamais.
\omcode{code/interviews-3/05-operating-systems-concurrency-and-networks/python/iv_arb.py}{1}{26}{Fusion de deux lignes.\protect\hfil\protect\break}\label{lst:iv3:os:arb}

La première copie l'emporte, si bien que le flux fusionné est au moins aussi rapide que la ligne la plus rapide, paquet par paquet, ce qui justifie d'en avoir deux. Chaque ligne est en général dans l'ordre~: un numéro manquant signifie donc surtout que cette ligne l'a perdu, et la copie de l'autre ligne arrive peu après.

\texttt{hold} fixe le compromis. Trop petit, et une copie simplement en retard sur la ligne la plus lente est déclarée perdue, ce qui déclenche une reprise inutile. Trop grand, et chaque perte réelle bloque toutes les données suivantes jusqu'à ce qu'elle soit reconnue. Il faut le choisir d'après l'écart mesuré entre les deux lignes~: combien de paquets peuvent arriver sur la ligne la plus rapide pendant le retard normal de la plus lente, aux débits de pointe. Un gestionnaire réel ajoute une minuterie, car pendant un marché calme aucun paquet ultérieur n'arrive pour déclencher le contrôle des trous. Cette version n'a pas de minuterie~: un trou en fin de rafale reste donc en attente.

Le trou déclaré doit aller quelque part~: une demande de retransmission, ou un instantané fourni par un service de reprise. Les copies tardives qui arrivent après le saut sont alors rejetées, car elles sont plus anciennes que \texttt{next}.

Le test fusionne deux lignes à pertes indépendantes, tirées au hasard, sur 300 cas. La sortie doit être croissante, sans doublon, et tous les numéros perdus sur les deux lignes et dépassés par la sortie doivent être signalés comme trous. Avec un grand \texttt{hold}, une perte sur une seule ligne ne devient jamais un trou.

\iqwrong{«~Transmettre ce qui arrive, et laisser le carnet d'ordres se débrouiller.~» Un carnet appliqué dans le désordre donne de faux prix, si bien qu'un seul flux ordonné doit sortir de cette couche.}
\end{solution}

\begin{solution}{iq:iv3:operating-systems-concurrency-and-networks:38}
Chaque connexion dispose d'un tampon de lecture, géré par son \texttt{Deframer}, et d'un tampon d'écriture, \texttt{out}. Les réponses en attente vont dans \texttt{out}. Les \cref{lst:iv3:os:loop1,lst:iv3:os:loop2} ne demandent \texttt{EPOLLOUT} que tant que \texttt{out} contient des données, car une socket inactive est presque toujours disponible en écriture, et la boucle se réveillerait sans arrêt. Le serveur ne se bloque jamais sur un seul client.
\omcode{code/interviews-3/05-operating-systems-concurrency-and-networks/python/iv_event_loop.py}{7}{38}{La boucle d'événements.\protect\hfil\protect\break}\label{lst:iv3:os:loop1}
\enlargethispage{-2\baselineskip}%
\omcode{code/interviews-3/05-operating-systems-concurrency-and-networks/python/iv_event_loop.py}{41}{68}{Lecture et envoi pour une connexion.\protect\hfil\protect\break}\label{lst:iv3:os:loop2}

La boucle utilise le mode par défaut, déclenché par niveau, si bien qu'un seul \texttt{recv} par événement suffit, et les données restantes donnent lieu à un nouvel événement. Toutes les sockets sont non bloquantes, ce qui compte surtout pour \texttt{send}~: sur une socket bloquante, un client qui cesse de lire arrêterait le serveur pour tous, dès que son tampon serait plein.

\texttt{send} peut n'accepter qu'une partie des données, donc \texttt{flush} retire ce qui a été accepté et conserve le reste. Quand \texttt{recv} renvoie une valeur vide, le pair a fini d'émettre, mais il peut encore lire, comme dans le cas d'une demi-fermeture. Le serveur cesse donc de lire, envoie ce qui se trouve dans \texttt{out}, puis ferme. Si le pair est réellement parti, un \texttt{recv} ou un \texttt{send} lève une erreur, et la connexion est abandonnée avec ses réponses en attente, qui n'ont plus de destination. Une longueur erronée venant d'un client ne ferme que ce client.

Dans le test, 8 clients envoient chacun 300 messages de 4000 octets au plus, découpés au hasard, puis effectuent une demi-fermeture, et lisent par morceaux de 1000 octets dans un thread séparé. Les sockets ont de petits tampons, si bien que le serveur doit mettre les réponses en file d'attente, et le test vérifie qu'il l'a fait. Chaque client récupère exactement ses messages, dans l'ordre, et le serveur s'arrête dès que tous ont fermé. Un second test fait fermer un client sans qu'il ait lu, et un autre envoyer une longueur supérieure à la limite, tandis qu'un troisième client est toujours servi.

Un client lent coûte peu de CPU aux autres, mais ce listing a un vrai défaut~: \texttt{out} n'a aucune limite, si bien qu'un client qui écrit sans jamais lire fait croître la mémoire du serveur sans borne. Un serveur réel la limite, en cessant de lire ce client dès que \texttt{out} est grand, ce qui remonte par TCP jusqu'à l'émetteur, ou en fermant la connexion. De plus, un message unique très volumineux retarde les autres, puisque la boucle n'exécute qu'un seul thread. Les gestionnaires ne doivent ni bloquer ni durer longtemps.

\iqwrong{«~Enregistrer \texttt{EPOLLOUT} pour toutes les sockets dès le départ.~» Avec le mode déclenché par niveau, la boucle tourne à pleine charge du CPU.}
\end{solution}

\begin{solution}{iq:iv3:operating-systems-concurrency-and-networks:39}
Le \Cref{lst:iv3:os:reliable} en est un exemple. L'émetteur renvoie chaque paquet non acquitté de sa fenêtre, le récepteur retient les paquets en avance et livre ceux qui se suivent. Il doit acquitter un doublon, car un doublon signifie que son ACK précédent a pu se perdre. Sinon l'émetteur renverrait sans fin.
\omcode{code/interviews-3/05-operating-systems-concurrency-and-networks/python/iv_reliable.py}{6}{43}{Répétition sélective, par manches.\protect\hfil\protect\break}\label{lst:iv3:os:reliable}

La fenêtre limite la quantité de données en transit, donc la quantité que le récepteur peut avoir à retenir et ce qu'une perte fait perdre. Avec $w = 1$, on a un protocole stop-and-wait, et le débit vaut un paquet par aller-retour, ce que le test vérifie à travers la fenêtre la plus large utilisée.

Pourquoi dans l'ordre et exactement une fois~? \texttt{delivered} ne croît que dans la boucle \texttt{while}, qui ajoute \texttt{next} puis l'incrémente, de sorte que les numéros livrés sont $0, 1, 2, \dots$, sans saut ni répétition. Un paquet de numéro inférieur à \texttt{next} n'est jamais stocké, si bien qu'un doublon ne peut pas ressortir. Pourquoi l'algorithme se termine-t-il~? Tant que le taux de perte est inférieur à 1, chaque paquet finit par passer avec son ACK, avec probabilité 1, si bien que la base continue d'avancer. L'ensemble \texttt{acked} de l'émetteur ne contient que des paquets que le récepteur a réellement reçus. Rien n'est donc abandonné à cause d'un faux acquittement.

Le test livre les 200 paquets à des taux de perte de 0, 0.1, 0.5 et 0.8, pour plusieurs graines aléatoires. Il vérifie la fenêtre, qu'en l'absence de perte il y a exactement 200 transmissions, qu'une perte plus forte coûte davantage, et qu'un taux de perte de 1 lève \texttt{TimeoutError} au lieu de bloquer. Les doublons et les réordonnancements sont prévus dans le modèle de liaison.

TCP a le même fond~: une fenêtre à acquittements cumulatifs, et des acquittements sélectifs en option. Ce que le modèle omet, c'est le temps réel~: les minuteries de retransmission, l'estimation du temps d'aller-retour, et une fenêtre de congestion qui se réduit à chaque perte. Le contrôle de flux est lui aussi absent~: ce récepteur a un tampon illimité. Dans un réseau réel, envoyer une fenêtre pleine à chaque manche, alors que la liaison est à l'origine de la perte, aggrave les choses.

\iqwrong{«~Numéroter les paquets, et retransmettre quand le récepteur le demande.~» Si la demande elle-même se perd, il ne se passe rien. Un des deux côtés a besoin d'une minuterie.}
\end{solution}

\begin{solution}{iq:iv3:operating-systems-concurrency-and-networks:40}
Le \Cref{lst:iv3:os:shmring} est le tampon circulaire. Les indices doivent être sans verrou, car un verrou à l'intérieur de \texttt{std::atomic} serait privé à un seul processus, ou resterait bloqué si son détenteur mourait. Le \texttt{static\_assert} sur \texttt{is\_always\_lock\_free} vérifie cela à la compilation. Si le producteur meurt, le consommateur voit tout ce qui a été publié, chaque case entière, et jamais une case à moitié écrite.
\omcode{code/interviews-3/05-operating-systems-concurrency-and-networks/cpp/shm_ring.hpp}{6}{27}{Un tampon circulaire en mémoire partagée.\protect\hfil\protect\break}\label{lst:iv3:os:shmring}

La case est écrite d'abord, et \texttt{head} n'avance qu'ensuite, par une écriture avec relâchement (release store). Ainsi, un arrêt forcé à n'importe quelle instruction laisse l'un de deux états~: \texttt{head} inchangé, et le contenu de la case simplement non vu, ou \texttt{head} avancé avec la case complète. Une écriture dans \texttt{head} est une seule étape atomique, sans état intermédiaire. Le consommateur ne peut pas savoir à partir du seul tampon que le producteur est mort, puisqu'un producteur mort et un producteur inactif se ressemblent, et il lui faut donc \texttt{waitpid}, un heartbeat ou un mécanisme similaire.

Le pilote fait pousser par un fils la suite 0, 1, 2 et ainsi de suite, jusqu'à ce que le parent le tue avec \texttt{SIGKILL}, à un moment qu'il n'a pas choisi. Le parent a alors vérifié 200\,000 valeurs, vide ce qui reste, et constate que la suite est ininterrompue. Il affiche \texttt{0 9 1}~: aucune erreur, signal 9, et assez de lectures.

Un mutex dans cette mémoire se comporterait différemment. Un processus tué en le détenant le laisse verrouillé, sauf s'il est robuste, et alors les données peuvent être à moitié mises à jour, si bien qu'il faut un plan de reprise. C'est pour cette raison qu'une voie de communication en mémoire partagée entre processus est souvent construite sans verrou.

ThreadSanitizer ne voit qu'un seul processus, si bien que le test exécute le même en-tête entre deux threads sous cet outil, et ne trouve aucune course. Avec les opérations store et load en relâchement, il signale une course aux données (\texttt{data race}), ce qui montre que le test sait détecter un mauvais ordre. Les ordres sont les mêmes, que ce soient des threads ou des processus qui les exécutent. L'architecture x86-64 masque ce type de défaut lors des exécutions réelles, et nous ne disposions pas ici de machine à ordre faible. L'alignement sur la ligne de cache empêche les deux indices de partager une même ligne.

\iqwrong{«~Après un plantage, le consommateur peut lire un message à moitié écrit.~» Seulement si \texttt{head} avance en premier. L'ordre des deux écritures est tout le principe de la conception.}
\end{solution}
